\SGASubsection{B.1}{The Frobenius morphism}
\SGARefTarget{sga3:xvii:appendix:B:subsection:1}

For every \(S\)-prescheme \(X\) and every integer \(n>0\), let
\[
  P^n:X\longrightarrow X
\]
denote the \(\mathbb F_p\)-endomorphism corresponding to raising
sections of \(\mathcal O_X\) to the \(p^n\)-th power, and let
\(X^{(n)}\) denote the \(S\)-prescheme obtained from \(X\) by base
change along \(P^n:S\to S\).  There is then a unique
\(S\)-morphism
\[
  F^n:X\longrightarrow X^{(n)}
\]
that makes the following diagram commutative:
\[
\begin{tikzcd}[column sep=3.7em,row sep=3.0em]
  {} & {} & X
    \arrow[dll,"P^n"']
    \arrow[dl,"F^n"]
    \arrow[ddl] \\
  X \arrow[d] &
  X^{(n)} \arrow[l] \arrow[d] &
  {} \\
  S &
  S \arrow[l,"P^n"] &
  {}
\end{tikzcd}
\]

It is clear that \(F^n\) identifies with the ``\(n\)-th iterate''
of \(F^1=F\), called the Frobenius morphism of \(X/S\).

If \(G\) is an \(S\)-group prescheme, then \(G^{(n)}\) is an
\(S\)-group prescheme and
\[
  F^n:G\longrightarrow G^{(n)}
\]
is an \(S\)-group morphism.  Its kernel
\({}_{F^n}(G)\) is a characteristic subgroup prescheme of \(G\)
(that is, it is stable under the functor
\(\operatorname{Aut}_{S\text{-}\mathrm{gr}}(G)\)), and is radicial
over \(S\).  If \(G\) is a radicial \(S\)-group prescheme, we say
that \(G\) has height at most \(h\) if
\({}_{F^h}(G)=G\).

\SGASubsection{B.2}{Groups and restricted
\texorpdfstring{\(p\)-Lie algebras}{p-Lie algebras}}
\SGARefTarget{sga3:xvii:appendix:B:subsection:2}

If \(G\) is an \(S\)-group prescheme, then
\(\operatorname{Lie}(G)\) (Exposé II) is naturally a restricted
\(p\)-Lie algebra (Exposé VII\textsubscript{A}, §§5 and 6).  In
particular, one has the following result (cf. Exposé
VII\textsubscript{A}).

\SGAStatement{\SGAInlineRefTarget{sga3:xvii:proposition:b-2-1}Proposition B.2.1}{%
\begin{enumerate}
\item[(i)] \(\operatorname{Lie}(\boldsymbol\alpha_p)_S
=\operatorname{Lie}(\mathbb G_a)_S\)
  is a sheaf of \(\mathcal O_S\)-modules, free of rank \(1\) over
  \(\mathcal O_S\), generated by an element \(X\) such that
  \(X^{(p)}=0\).
\item[(ii)] \(\operatorname{Lie}(\boldsymbol\mu_p)_S
  =\operatorname{Lie}(\mathbb G_m)_S\) is a sheaf of
  \(\mathcal O_S\)-modules, free of rank \(1\) over
  \(\mathcal O_S\), with a canonical basis \(X\) such that
  \(X^{(p)}=X\).
\end{enumerate}}
\SGARefTarget{sga3:xvii:proposition:B:2:1}

We now recall the fundamental result proved in Exposé
VII\textsubscript{A}, §7.

\SGAStatement{\SGAInlineRefTarget{sga3:xvii:theorem:b-2-2}Theorem B.2.2}{Suppose that \(S\) is affine, with
ring \(A\).  Then the functor
\[
  G\longmapsto\operatorname{Lie}G
\]
defines an equivalence between, on the one hand, the category of
\(S\)-group preschemes \(G\) that are of finite presentation and
flat over \(S\), of height \(1\), and whose Lie algebra is locally
free over \(\mathcal O_S\), and, on the other hand, the category
of restricted \(p\)-\(A\)-Lie algebras that are locally free of
finite rank.

Moreover, if \(G\) is as above and \(H\) is any \(S\)-group
prescheme of finite presentation, the canonical morphism
\[
  \operatorname{Hom}_{S\text{-}\mathrm{gr}}(G,H)
  \longrightarrow
  \operatorname{Hom}_{p\text{-}A\text{-}\mathrm{Lie}}
    \bigl(\operatorname{Lie}G(S),\operatorname{Lie}H(S)\bigr)
\]
is an isomorphism.}
\SGARefTarget{sga3:xvii:theorem:B:2:2}

\SGANextSectionTarget{sga3:xvii:subsection:b-3}\SGASubsection{B.3}{Radicial groups and smooth groups}
\SGARefTarget{sga3:xvii:appendix:B:subsection:3}

We now suppose that \(S\) is the spectrum of a field \(k\) of
characteristic \(p\).  Recall that Exposé
VI\textsubscript{A}, §5 proves that, if \(G\) is an algebraic
\(k\)-group and \(H\) an algebraic subgroup of \(G\), then the
sheaf \(G/H\) (for the fpqc topology) is representable.  We also
recall (VII\textsubscript{A}, \SGARefLink{sga3:VIIA:8.3}{8.3}):

\SGAStatement{\SGAInlineRefTarget{sga3:xvii:proposition:b-3-1}Proposition B.3.1}{Let \(G\) be an algebraic
\(k\)-group.  There exists an integer \(m\) such that, for every
\(n\geq m\), the algebraic group
\[
  G/{}_{F^n}(G)
\]
is smooth over \(k\).}
\SGARefTarget{sga3:xvii:proposition:B:3:1}

\SGAStatement{Proposition B.3.2}{Consider an exact sequence of
algebraic \(k\)-groups
\[
  1\longrightarrow G'\longrightarrow G
    \xrightarrow{u}G''\longrightarrow1
\]
and the following assertions:
\begin{enumerate}
\item[(i)] \SGARefTarget{sga3:xvii:proposition:B:3:2:item:i}
  the morphism \(u\) is smooth;
\item[(ii)] \SGARefTarget{sga3:xvii:proposition:B:3:2:item:ii}
  \(G'\) is smooth over \(k\);
\item[(iii)] \SGARefTarget{sga3:xvii:proposition:B:3:2:item:iii}
  for every integer \(n>0\), the sequence
  \[
    1\longrightarrow{}_{F^n}(G')
      \longrightarrow{}_{F^n}(G)
      \longrightarrow{}_{F^n}(G'')
      \longrightarrow1
  \]
  is exact;
\item[(iv)] \SGARefTarget{sga3:xvii:proposition:B:3:2:item:iv}
  the morphism \({}_FG\to{}_FG''\) is an epimorphism;
\item[(v)] \SGARefTarget{sga3:xvii:proposition:B:3:2:item:v}
  the morphism
  \[
    \bigl(\operatorname{Lie}G\bigr)(k)
      \longrightarrow\bigl(\operatorname{Lie}G''\bigr)(k)
  \]
  is surjective.
\end{enumerate}
Then
\[
  \text{\rm(i)}\Longleftrightarrow\text{\rm(ii)}
  \Longrightarrow\text{\rm(iii)}
  \Longrightarrow\text{\rm(iv)}
  \Longleftrightarrow\text{\rm(v)}.
\]
Moreover, if \(G\) is smooth, all five assertions are equivalent.}
\SGARefTarget{sga3:xvii:proposition:B:3:2}

The equivalence \SGARefLink{sga3:xvii:proposition:B:3:2:item:i}{(i)}\(\Longleftrightarrow\)\SGARefLink{sga3:xvii:proposition:B:3:2:item:ii}{(ii)} follows from
Exposé VI\textsubscript{B}, \SGARefLink{concl:sga3-VIB-9.2-vii}{9.2(vii)}.

For \SGARefLink{sga3:xvii:proposition:B:3:2:item:ii}{(ii)}\(\Longrightarrow\)\SGARefLink{sga3:xvii:proposition:B:3:2:item:iii}{(iii)}, if \(G'\) is smooth, then
\[
  F^n:G'\longrightarrow G'^{(n)}
\]
is an epimorphism, and \SGARefLink{sga3:xvii:proposition:B:3:2:item:iii}{(iii)} follows from the snake diagram.
The implication \SGARefLink{sga3:xvii:proposition:B:3:2:item:iii}{(iii)}\(\Longrightarrow\)\SGARefLink{sga3:xvii:proposition:B:3:2:item:iv}{(iv)} is clear, while
\SGARefLink{sga3:xvii:proposition:B:3:2:item:iv}{(iv)}\(\Longleftrightarrow\)\SGARefLink{sga3:xvii:proposition:B:3:2:item:v}{(v)} follows from \SGARefLink{sga3:xvii:theorem:b-2-2}{Theorem B.2.2}.
Finally, suppose that \(G\) is smooth and that \SGARefLink{sga3:xvii:proposition:B:3:2:item:v}{(v)} holds.  Then
\(G''\) is smooth, and \SGARefLink{sga3:xvii:proposition:B:3:2:item:v}{(v)} implies
\[
  \dim G'=\dim_k\bigl(\operatorname{Lie}G'\bigr)(k),
\]
so \(G'\) is smooth.

\SGASection{C}{Appendix III. Remarks and Supplements Concerning
Exposés XV, XVI, and XVII}
\SGARefTarget{sga3:xvii:appendix:C}

\SGANextSectionTarget{sga3:xvii:subsection:c-1}\SGASubsection{C.1}{}
\SGARefTarget{sga3:xvii:appendix:C:subsection:1}

It may be that \SGARefLink{sga3:xvii:proposition:1:2}{Propositions 1.2} and \SGARefLink{sga3:xv:proposition:1-2-bis}{1.2 bis} of Exposé XV remain
valid if the hypothesis that \(H_0\) is smooth over \(S_0\) is
removed.  This is the case, in particular, if \(G\) is finite,
flat, and commutative.
